\section{Suplemento sobre la cohomología de haces}
\label{section:0.12}

\subsection{Cohomología de haces de módulos sobre espacios anillados}
\label{subsection:0.12.1}

\begin{env}[12.1.1]
\label{0.12.1.1}
\oldpage[0\textsubscript{III}]{49}
Sea $(X,\sh{O}_X)$ un espacio anillado.
Recordemos que, para todo $\sh{O}_X$-módulo $\sh{F}$, se define la
cohomología $\HH^\bullet(X,\sh{F})$, que es un funtor cohomológico universal
(T, 2.2) de la categoría $\cat{C}(X)$ de
$\sh{O}_X$-módulos en la categoría de grupos abelianos; es el funtor derivado
del funtor exacto por la izquierda
\[
  \sh{F}\longmapsto\Gamma(X,\sh{F}).
\]
El funtor $\HH^\bullet(X,\sh{F})$ es isomorfo
\oldpage[0\textsubscript{III}]{50}
a la restricción a la categoría $\cat{C}(X)$ del funtor cohomológico
definido del mismo modo sobre la categoría de haces de grupos abelianos
en $X$ (G, II, 7.2.1).
\end{env}

\begin{env}[12.1.2]
\label{0.12.1.2}
Pongamos $A=\Gamma(X,\sh{O}_X)$.
Como todo elemento de $A$ define un endomorfismo del grupo abeliano
$\Gamma(X,\sh{F})$, define por funtorialidad un endomorfismo del
$\delta$-funtor $\HH^\bullet(X,\sh{F})$.
Estos endomorfismos dotan a cada $\HH^p(X,\sh{F})$ de una
estructura de $A$-módulo, y el operador de conexión $\partial$ es
$A$-lineal.
Además, para cualesquiera dos enteros no negativos $p,q$ y cualesquiera dos
$\sh{O}_X$-módulos $\sh{F},\sh{G}$, existe un homomorfismo de
$A$-módulos, llamado \emph{producto cup},
\[
\label{0.12.1.2.1}
\tag{12.1.2.1}
  \HH^p(X,\sh{F})\otimes_A\HH^q(X,\sh{G})
  \longrightarrow
  \HH^{p+q}\bigl(X,\sh{F}\otimes_{\sh{O}_X}\sh{G}\bigr)
\]
(G, II, 6.6).
Estos homomorfismos hacen de la suma directa $S$ de los
$\HH^p(X,\sh{O}_X)$, para $p\geq0$, una $A$-álgebra graduada
anticomutativa, y de la suma directa de los $\HH^p(X,\sh{F})$ un
$S$-módulo graduado.

Para todo recubrimiento abierto $\mathfrak{U}$ de $X$, denotaremos siempre
por $\check{\mathrm{C}}^\bullet(\mathfrak{U},\sh{F})$ (contrariamente a
(G, II, 5.1)) el complejo de cocadenas alternadas del nervio de
$\mathfrak{U}$ con valores en el sistema de coeficientes
$\Gamma(U_\alpha,\sh{F})$.
Está claro que
$\check{\mathrm{C}}^\bullet(\mathfrak{U},\sh{F})$ es un
$A$-módulo graduado, de modo que los grupos de cohomología
$\check{\HH}^p(\mathfrak{U},\sh{F})$ de este complejo están dotados de
una estructura de $A$-módulo.
Además, las aplicaciones canónicas
\[
  \check{\HH}^p(\mathfrak{U},\sh{F})
  \longrightarrow
  \HH^p(X,\sh{F})
\]
(G, II, 5.4) son $A$-lineales.
\end{env}

\begin{env}[12.1.3]
\label{0.12.1.3}
Sea $(X',\sh{O}_{X'})$ un segundo espacio anillado, y sea
$f=(\psi,\theta)$ un morfismo de $X'$ en $X$.
Pongamos $A'=\Gamma(X',\sh{O}_{X'})$; el $\psi$-morfismo $\theta$
define canónicamente un homomorfismo de anillos $A\to A'$.
Sea $\sh{F}$ un $\sh{O}_X$-módulo y $\sh{F}'$ un
$\sh{O}_{X'}$-módulo.
Para todo $f$-morfismo $u:\sh{F}\to\sh{F}'$
\sref[0\textsubscript{I}]{0.4.4.1}, veremos que se puede definir,
para todo $p\geq0$, un dihomomorfismo
\[
\label{0.12.1.3.1}
\tag{12.1.3.1}
  u_p:\HH^p(X,\sh{F})\longrightarrow\HH^p(X',\sh{F}').
\]

En efecto, como $\psi^*$ es exacto en la categoría de haces de grupos
abelianos sobre $X$, el funtor
\[
  \sh{F}\longmapsto\HH^\bullet(X',\psi^*(\sh{F}))
\]
es un $\delta$-funtor sobre dicha categoría, y se sabe que existe un
homomorfismo canónico de $\delta$-funtores
\[
\label{0.12.1.3.2}
\tag{12.1.3.2}
  \HH^\bullet(X,\sh{F})
  \longrightarrow
  \HH^\bullet(X',\psi^*(\sh{F}))
\]
determinado de manera única por la condición de que, en grado $0$, se reduce
al homomorfismo canónico
\[
  \Gamma(X,\sh{F})
  \longrightarrow
  \Gamma(X',\psi^*(\sh{F}))
\]
(T, 3.2.2).
Además, todo elemento de $A$ determina un endomorfismo $\mu$ de
$\Gamma(X,\sh{F})$ y un endomorfismo $\mu'$ de
$\Gamma(X',\psi^*(\sh{F}))$ tal que el diagrama
\[
\label{0.12.1.3.3}
\tag{12.1.3.3}
  \xymatrix{
    \Gamma(X,\sh{F})\ar[r]\ar[d]_\mu
      & \Gamma(X',\psi^*(\sh{F}))\ar[d]^{\mu'}\\
    \Gamma(X,\sh{F})\ar[r]
      & \Gamma(X',\psi^*(\sh{F}))
  }
\]
es conmutativo.
\oldpage[0\textsubscript{III}]{51}
Por la propiedad de unicidad para prolongar morfismos de funtores
cohomológicos universales (T, 2.2), se obtienen prolongaciones únicas de
$\mu$ y $\mu'$ a la cohomología que hacen conmutativos los diagramas
análogos a \sref{0.12.1.3.3}.
Así, \sref{0.12.1.3.2} es un homomorfismo de $A$-módulos.

Observemos ahora que
\[
  f^*(\sh{F})
  =
  \psi^*(\sh{F})
  \otimes_{\psi^*(\sh{O}_X)}
  \sh{O}_{X'},
\]
y, por tanto, que existe un dihomomorfismo canónico
\[
  \psi^*(\sh{F})\longrightarrow f^*(\sh{F})
\]
del $\psi^*(\sh{O}_X)$-módulo $\psi^*(\sh{F})$ en el
$\sh{O}_{X'}$-módulo $f^*(\sh{F})$.
Por funtorialidad se obtiene, pues, un dihomomorfismo funtorial
\[
\label{0.12.1.3.4}
\tag{12.1.3.4}
  \HH^p(X',\psi^*(\sh{F}))
  \longrightarrow
  \HH^p(X',f^*(\sh{F})),
\]
siendo $A$ y $A'$ los anillos correspondientes.
Componiendo este dihomomorfismo con \sref{0.12.1.3.2}, se obtiene un
dihomomorfismo canónico, funtorial en $\sh{F}$,
\[
\label{0.12.1.3.5}
\tag{12.1.3.5}
  \theta_p:\HH^p(X,\sh{F})
  \longrightarrow
  \HH^p(X',f^*(\sh{F})).
\]

Finalmente, por funtorialidad, el homomorfismo
$u^\sharp:f^*(\sh{F})\to\sh{F}'$ da un homomorfismo de
$A'$-módulos
\[
  \HH^p(X',f^*(\sh{F}))
  \longrightarrow
  \HH^p(X',\sh{F}'),
\]
que, compuesto con \sref{0.12.1.3.5}, da
\sref{0.12.1.3.1}.

Sea $f'=(\psi',\theta'):X''\to X'$ un segundo morfismo de espacios
anillados, y sea $f''=f\circ f'$ el morfismo compuesto.
Teniendo en cuenta la compatibilidad del funtor $\psi'^*$ con los
productos tensoriales \sref[0\textsubscript{I}]{0.4.3.3}, se comprueba
inmediatamente que el compuesto del dihomomorfismo
\[
  \HH^p(X',f^*(\sh{F}))
  \longrightarrow
  \HH^p\bigl(X'',f'^*(f^*(\sh{F}))\bigr)
\]
con \sref{0.12.1.3.5} es el dihomomorfismo correspondiente
\[
  \HH^p(X,\sh{F})
  \longrightarrow
  \HH^p(X'',f''^*(\sh{F})).
\]
\end{env}

\begin{env}[12.1.4]
\label{0.12.1.4}
Puede darse una definición directa del homomorfismo
\sref{0.12.1.3.2} como sigue.
Tomemos una resolución inyectiva
$\sh{L}^\bullet=(\sh{L}^i)$ de $\sh{F}$ constituida por haces de
grupos abelianos sobre $X$.
Como el funtor $\psi^*$ es exacto,
$\psi^*(\sh{L}^\bullet)$ es una resolución de $\psi^*(\sh{F})$
constituida por haces sobre $X'$.
Si $\sh{L}'^\bullet=(\sh{L}'^i)$ es una resolución inyectiva de
$\psi^*(\sh{F})$ en la categoría de haces de grupos abelianos sobre
$X'$, existe un morfismo
\[
  \psi^*(\sh{L}^\bullet)\longrightarrow\sh{L}'^\bullet
\]
de complejos de haces de grupos abelianos, compatible con los
aumentos (M, V, 1.1.a), y determinado salvo homotopía.
Se obtienen así homomorfismos
\[
  \Gamma(X,\sh{L}^\bullet)
  \longrightarrow
  \Gamma(X',\psi^*(\sh{L}^\bullet))
  \longrightarrow
  \Gamma(X',\sh{L}'^\bullet)
\]
de complejos de grupos abelianos.
Su compuesto, al pasar a la cohomología, da un morfismo de
$\delta$-funtores
\[
  \HH^\bullet(X,\sh{F})
  \longrightarrow
  \HH^\bullet(X',\psi^*(\sh{F}));
\]
como coincide con \sref{0.12.1.3.2} en grado $0$, es idéntico a él
(T, 2.2).

Consideremos ahora un recubrimiento abierto
$\mathfrak{U}=(U_\alpha)$ de $X$, y sea
$\mathfrak{U}'=(f^{-1}(U_\alpha))$ su recubrimiento abierto imagen inversa
de $X'$.
Para todo subconjunto abierto $V$ de $X$, los homomorfismos canónicos
\[
  \Gamma(V,\sh{F})
  \longrightarrow
  \Gamma(f^{-1}(V),f^*(\sh{F}))
\]
definen inmediatamente (cf.\ G, II, 5.1) un homomorfismo de complejos
\[
  \check{\mathrm{C}}^\bullet(\mathfrak{U},\sh{F})
  \longrightarrow
  \check{\mathrm{C}}^\bullet(\mathfrak{U}',f^*(\sh{F})),
\]
y, por tanto, homomorfismos canónicos
\[
\label{0.12.1.4.1}
\tag{12.1.4.1}
  \theta_p:
  \check{\HH}^p(\mathfrak{U},\sh{F})
  \longrightarrow
  \check{\HH}^p(\mathfrak{U}',f^*(\sh{F})).
\]

\oldpage[0\textsubscript{III}]{52}
Además, existen diagramas conmutativos
\[
\label{0.12.1.4.2}
\tag{12.1.4.2}
  \xymatrix{
    \check{\HH}^p(\mathfrak{U},\sh{F})
      \ar[r]^-{\theta_p}\ar[d]
      & \check{\HH}^p(\mathfrak{U}',f^*(\sh{F}))\ar[d]\\
    \HH^p(X,\sh{F})
      \ar[r]_-{\theta_p}
      & \HH^p(X',f^*(\sh{F})),
  }
\]
donde las flechas verticales son los homomorfismos canónicos de
(G, II, 5.2).
Para demostrar la conmutatividad de \sref{0.12.1.4.2}, consideremos el
complejo de haces de cocadenas alternadas de $\sh{F}$ relativo a
$\mathfrak{U}$,
$\sh{C}^\bullet(\mathfrak{U},\sh{F})$, para el cual
\[
  \Gamma\bigl(X,\sh{C}^\bullet(\mathfrak{U},\sh{F})\bigr)
  =
  \check{\mathrm{C}}^\bullet(\mathfrak{U},\sh{F})
\]
(G, II, 5.2).
Los homomorfismos canónicos
\[
  \Gamma(V,\sh{F})
  \longrightarrow
  \Gamma(f^{-1}(V),\psi^*(\sh{F}))
\]
definen entonces un $\psi$-morfismo
\[
  \sh{C}^\bullet(\mathfrak{U},\sh{F})
  \longrightarrow
  \sh{C}^\bullet(\mathfrak{U}',\psi^*(\sh{F})),
\]
y, con la notación precedente, existe un diagrama conmutativo
\[
  \xymatrix{
    \Gamma\bigl(X,\sh{C}^\bullet(\mathfrak{U},\sh{F})\bigr)
      \ar[r]\ar[d]
      & \Gamma\bigl(X',
          \sh{C}^\bullet(\mathfrak{U}',\psi^*(\sh{F}))\bigr)\ar[d]\\
    \Gamma(X,\sh{L}^\bullet)\ar[r]
      & \Gamma(X',\psi^*(\sh{L}^\bullet))\ar[r]
      & \Gamma(X',\sh{L}'^\bullet).
  }
\]
Al pasar a la cohomología se obtienen diagramas conmutativos
\[
  \xymatrix{
    \check{\HH}^p(\mathfrak{U},\sh{F})\ar[r]\ar[d]
      & \check{\HH}^p(\mathfrak{U}',\psi^*(\sh{F}))\ar[d]\\
    \HH^p(X,\sh{F})\ar[r]
      & \HH^p(X',\psi^*(\sh{F})),
  }
\]
donde las flechas verticales son los homomorfismos canónicos de
(G, II, 5.2).
Solo queda combinar estos diagramas con los diagramas conmutativos
\[
  \xymatrix{
    \check{\HH}^p(\mathfrak{U}',\psi^*(\sh{F}))\ar[r]\ar[d]
      & \check{\HH}^p(\mathfrak{U}',f^*(\sh{F}))\ar[d]\\
    \HH^p(X',\psi^*(\sh{F}))\ar[r]
      & \HH^p(X',f^*(\sh{F})),
  }
\]
\oldpage[0\textsubscript{III}]{53}
que provienen del homomorfismo
$\psi^*(\sh{F})\to f^*(\sh{F})$ y de la funtorialidad de los
homomorfismos canónicos de (G, II, 5.2), para obtener la conmutatividad
de \sref{0.12.1.4.2}.

Observemos que si
$A=\Gamma(X,\sh{O}_X)$ y $A'=\Gamma(X',\sh{O}_{X'})$, entonces
\sref{0.12.1.4.1} es un dihomomorfismo de módulos correspondiente a
los anillos $A$ y $A'$.
Existe para \sref{0.12.1.4.1} una propiedad de transitividad respecto
de un compuesto de dos morfismos, análoga a la de
\sref{0.12.1.3.5}.
Finalmente, en las definiciones precedentes, en lugar de una resolución
inyectiva $\sh{L}^\bullet$ de $\sh{F}$, se podría haber partido igualmente
de una resolución tal que
\[
  \HH^p(X,\sh{L}^i)=0
  \qquad\text{para todo $i$ y todo $p>0$}
\]
(G, II, 4.7.1).
\end{env}

\begin{env}[12.1.5]
\label{0.12.1.5}
Sean $\sh{F},\sh{G},\sh{H}$ tres $\sh{O}_X$-módulos, y consideremos
un $\sh{O}_X$-homomorfismo
\[
  u:\sh{F}\otimes_{\sh{O}_X}\sh{G}\longrightarrow\sh{H},
\]
que da, en cohomología, los homomorfismos
\[
\label{0.12.1.5.1}
\tag{12.1.5.1}
  \HH^p(X,\sh{F})\otimes_A\HH^q(X,\sh{G})
  \longrightarrow
  \HH^{p+q}(X,\sh{H})
\]
deducidos del producto cup \sref{0.12.1.2.1}.
Demostraremos que, bajo las hipótesis y con la notación de
\sref{0.12.1.3}, los diagramas
\[
\label{0.12.1.5.2}
\tag{12.1.5.2}
  \xymatrix@C=10pt{
    \HH^p(X,\sh{F})\otimes_A\HH^q(X,\sh{G})
      \ar[r]\ar[d]
      & \HH^{p+q}(X,\sh{H})\ar[d]\\
    \HH^p(X',f^*(\sh{F}))\otimes_{A'}
      \HH^q(X',f^*(\sh{G}))
      \ar[r]
      & \HH^{p+q}(X',f^*(\sh{H}))
  }
\]
son conmutativos, donde las flechas verticales provienen de los
homomorfismos canónicos \sref{0.12.1.3.5}.

Recordemos para ello que \sref{0.12.1.5.1} puede obtenerse a partir de
las resoluciones canónicas (G, II, 4.3)
$\sh{L}^\bullet,\sh{M}^\bullet,\sh{N}^\bullet$ de
$\sh{F},\sh{G},\sh{H}$ respectivamente, que están constituidas por
$\sh{O}_X$-módulos, y de la aplicación lineal
\[
  \sh{L}^\bullet\otimes_{\sh{O}_X}\sh{M}^\bullet
  \longrightarrow
  \sh{N}^\bullet
\]
de complejos de $\sh{O}_X$-módulos correspondiente a $u$.
Esta aplicación proporciona un homomorfismo de complejos de $A$-módulos
\[
  \Gamma(X,\sh{L}^\bullet)\otimes_A\Gamma(X,\sh{M}^\bullet)
  \longrightarrow
  \Gamma(X,\sh{N}^\bullet),
\]
y, por tanto, al pasar a la cohomología, homomorfismos
\[
  \HH^p\bigl(\Gamma(X,\sh{L}^\bullet)\bigr)
  \otimes_A
  \HH^q\bigl(\Gamma(X,\sh{M}^\bullet)\bigr)
  \longrightarrow
  \HH^{p+q}\bigl(\Gamma(X,\sh{N}^\bullet)\bigr)
\]
(G, II, 6.6).
Existe claramente un diagrama conmutativo
\[
\label{0.12.1.5.3}
\tag{12.1.5.3}
  \xymatrix@C=8pt{
    \Gamma(X,\sh{L}^\bullet)\otimes_A
      \Gamma(X,\sh{M}^\bullet)
      \ar[r]\ar[d]
      & \Gamma(X,\sh{N}^\bullet)\ar[d]\\
    \Gamma(X',\psi^*(\sh{L}^\bullet))
      \otimes_{\Gamma(X',\psi^*(\sh{O}_X))}
      \Gamma(X',\psi^*(\sh{M}^\bullet))
      \ar[r]
      & \Gamma(X',\psi^*(\sh{N}^\bullet)).
  }
\]
\oldpage[0\textsubscript{III}]{54}
Al pasar a la cohomología se obtienen diagramas conmutativos
\[
\label{0.12.1.5.4}
\tag{12.1.5.4}
  \xymatrix@C=7pt{
    \HH^p(X,\sh{F})\otimes_A\HH^q(X,\sh{G})
      \ar[r]\ar[d]
      & \HH^{p+q}(X,\sh{H})\ar[d]\\
    \HH^p\bigl(\Gamma(X',\psi^*(\sh{L}^\bullet))\bigr)
      \otimes_{\Gamma(X',\psi^*(\sh{O}_X))}
      \HH^q\bigl(\Gamma(X',\psi^*(\sh{M}^\bullet))\bigr)
      \ar[r]
      & \HH^{p+q}\bigl(\Gamma(X',\psi^*(\sh{N}^\bullet))\bigr).
  }
\]
Pero como
$\psi^*(\sh{L}^\bullet)$,
$\psi^*(\sh{M}^\bullet)$, y
$\psi^*(\sh{N}^\bullet)$
son resoluciones de
$\psi^*(\sh{F})$,
$\psi^*(\sh{G})$, y
$\psi^*(\sh{H})$ respectivamente, existe un diagrama conmutativo
(G, II, 6.6.1)
\[
\label{0.12.1.5.5}
\tag{12.1.5.5}
  \xymatrix@C=7pt{
    \HH^p\bigl(\Gamma(X',\psi^*(\sh{L}^\bullet))\bigr)
      \otimes_{\Gamma(X',\psi^*(\sh{O}_X))}
      \HH^q\bigl(\Gamma(X',\psi^*(\sh{M}^\bullet))\bigr)
      \ar[r]\ar[d]
      & \HH^{p+q}\bigl(\Gamma(X',\psi^*(\sh{N}^\bullet))\bigr)
        \ar[d]\\
    \HH^p(X',\psi^*(\sh{F}))
      \otimes_{\Gamma(X',\psi^*(\sh{O}_X))}
      \HH^q(X',\psi^*(\sh{G}))
      \ar[r]
      & \HH^{p+q}(X',\psi^*(\sh{H})).
  }
\]
Finalmente, por funtorialidad, existe un diagrama conmutativo
\[
\label{0.12.1.5.6}
\tag{12.1.5.6}
  \xymatrix@C=7pt{
    \HH^p(X',\psi^*(\sh{F}))
      \otimes_{\Gamma(X',\psi^*(\sh{O}_X))}
      \HH^q(X',\psi^*(\sh{G}))
      \ar[r]\ar[d]
      & \HH^{p+q}(X',\psi^*(\sh{H}))\ar[d]\\
    \HH^p(X',f^*(\sh{F}))\otimes_{A'}
      \HH^q(X',f^*(\sh{G}))
      \ar[r]
      & \HH^{p+q}(X',f^*(\sh{H})).
  }
\]
Combinando los tres diagramas \sref{0.12.1.5.4},
\sref{0.12.1.5.5} y \sref{0.12.1.5.6}, se obtiene el diagrama
conmutativo deseado \sref{0.12.1.5.2}.
\end{env}

\begin{remark}[12.1.6]
\label{0.12.1.6}
\oldpage[0\textsubscript{III}]{55}
Con la notación de \sref{0.12.1.3}, supongamos que existe un
diagrama conmutativo
\[
\label{0.12.1.6.1}
\tag{12.1.6.1}
  \xymatrix{
    0\ar[r] & \sh{F}\ar[r]^r\ar[d]_u
      & \sh{G}\ar[r]^s\ar[d]_v
      & \sh{H}\ar[r]\ar[d]_w & 0\\
    0\ar[r] & \sh{F}'\ar[r]_{r'}
      & \sh{G}'\ar[r]_{s'}
      & \sh{H}'\ar[r] & 0,
  }
\]
donde $r,s$ son homomorfismos de $\sh{O}_X$-módulos,
$r',s'$ son homomorfismos de $\sh{O}_{X'}$-módulos,
$u,v,w$ son $f$-morfismos, y las filas son exactas.
Se sigue que el diagrama
\[
\label{0.12.1.6.2}
\tag{12.1.6.2}
  \xymatrix@C=6pt{
    \cdots\ar[r]
      & \HH^p(X,\sh{F})\ar[r]\ar[d]_{u_p}
      & \HH^p(X,\sh{G})\ar[r]\ar[d]_{v_p}
      & \HH^p(X,\sh{H})\ar[r]^-\partial\ar[d]_{w_p}
      & \HH^{p+1}(X,\sh{F})\ar[r]\ar[d]_{u_{p+1}}
      & \cdots\\
    \cdots\ar[r]
      & \HH^p(X',\sh{F}')\ar[r]
      & \HH^p(X',\sh{G}')\ar[r]
      & \HH^p(X',\sh{H}')\ar[r]_-\partial
      & \HH^{p+1}(X',\sh{F}')\ar[r]
      & \cdots
  }
\]
es conmutativo.
En efecto, \sref{0.12.1.6.1} se factoriza como
\[
  \xymatrix@R=14pt{
    0\ar[r] & \sh{F}\ar[r]\ar[d]
      & \sh{G}\ar[r]\ar[d]
      & \sh{H}\ar[r]\ar[d] & 0\\
    0\ar[r] & \psi^*(\sh{F})\ar[r]\ar[d]
      & \psi^*(\sh{G})\ar[r]\ar[d]
      & \psi^*(\sh{H})\ar[r]\ar[d] & 0\\
    0\ar[r] & \sh{F}'\ar[r]
      & \sh{G}'\ar[r]
      & \sh{H}'\ar[r] & 0,
  }
\]
donde la fila intermedia es exacta
\sref[0\textsubscript{I}]{0.3.7.2}.
Basta entonces utilizar el hecho de que \sref{0.12.1.3.2} es un
homomorfismo de $\delta$-funtores y que los
$\HH^p(X',\sh{F}')$ forman un $\delta$-funtor en $\sh{F}'$.
\end{remark}

\begin{env}[12.1.7]
\label{0.12.1.7}
Con las hipótesis y la notación de \sref{0.12.1.3}, consideremos ahora el
caso
\[
  \sh{F}=f_*(\sh{F}')=\psi_*(\sh{F}').
\]
Veremos que el dihomomorfismo definido en \sref{0.12.1.3},
\[
  \HH^p(X,f_*(\sh{F}'))\longrightarrow\HH^p(X',\sh{F}'),
  \tag{12.1.7.1}
  \label{0.12.1.7.1}
\]
puede obtenerse (salvo un automorfismo de $\HH^p(X',\sh{F}')$) como un
homomorfismo de borde de una sucesión espectral para el funtor compuesto
\[
  \sh{F}'\leadsto\Gamma(X,\psi_*(\sh{F}'))
\]
(T, 2.4).
La descripción de \sref{0.12.1.7.1} dada en \sref{0.12.1.4} muestra
aquí que este homomorfismo puede obtenerse como sigue.
Tomemos resoluciones inyectivas $\sh{L}^\bullet$ y
${\sh{L}'}^\bullet$ de $\psi_*(\sh{F}')$ y $\sh{F}'$,
respectivamente, y después un homomorfismo de complejos
\[
  v:\psi^*(\sh{L}^\bullet)\longrightarrow{\sh{L}'}^\bullet
\]
``por encima'' del homomorfismo canónico
$\psi^*(\psi_*(\sh{F}'))\to\sh{F}'$.
\oldpage[0\textsubscript{III}]{56}
Observemos a continuación que
\[
  \Gamma(X',{\sh{L}'}^\bullet)
  =\Gamma(X,\psi_*({\sh{L}'}^\bullet))
\]
y que el homomorfismo compuesto
\[
  \Gamma(X,\sh{L}^\bullet)
  \longrightarrow\Gamma(X',\psi^*(\sh{L}^\bullet))
  \xrightarrow{\Gamma(v)}\Gamma(X',{\sh{L}'}^\bullet)
\]
no es otro que
\[
  \Gamma(v^\flat):
  \Gamma(X,\sh{L}^\bullet)
  \longrightarrow\Gamma(X,\psi_*({\sh{L}'}^\bullet))
  \tag{12.1.7.2}
  \label{0.12.1.7.2}
\]
\sref[0\textsubscript{I}]{0.3.7.1}; y \sref{0.12.1.7.1} se
obtiene pasando a la cohomología en \sref{0.12.1.7.2}.

Por otra parte, las sucesiones espectrales para el funtor compuesto
\[
  \sh{F}'\leadsto\Gamma(X,\psi_*(\sh{F}'))
\]
se obtienen tomando una resolución inyectiva de Cartan--Eilenberg
\[
  \sh{M}^{\bullet\bullet}=(\sh{M}^{ij})
\]
del complejo $\psi_*({\sh{L}'}^\bullet)$ en la categoría de
haces de grupos abelianos sobre $X$.
Las sucesiones espectrales en cuestión son las del bicomplejo
$\Gamma(X,\sh{M}^{\bullet\bullet})$; son birregulares, pues
$\sh{M}^{ij}=0$ para $i<0$ o $j<0$.
Ahora bien, la primera sucesión espectral de este bicomplejo degenera, pues
los haces $\psi_*({\sh{L}'}^i)$ son flácidos (G, II, 3.1.1), y
por tanto
\[
  \HH_\mathrm{II}^q
    \bigl(\Gamma(X,\sh{M}^{i,\bullet})\bigr)
  =\HH^q\bigl(\psi_*({\sh{L}'}^i)\bigr)=0
  \qquad(q>0)
\]
(G, II, 4.4.3).
Por consiguiente, existen homomorfismos de borde biyectivos
\sref{0.11.1.6}
\[
  {}'E_2^{i0}
  =\HH^i\!\left(
      \HH_\mathrm{II}^0
        \bigl(\Gamma(X,\sh{M}^{\bullet\bullet})\bigr)
    \right)
  \longrightarrow
  \HH^i\bigl(\Gamma(X,\sh{M}^{\bullet\bullet})\bigr),
  \tag{12.1.7.3}
  \label{0.12.1.7.3}
\]
y sabemos por \sref{0.11.3.4} que este homomorfismo proviene, al pasar
a la cohomología, del aumento
\[
  \Gamma(X,\psi_*({\sh{L}'}^\bullet))
  \longrightarrow\Gamma(X,\sh{M}^{\bullet\bullet}),
  \tag{12.1.7.4}
  \label{0.12.1.7.4}
\]
que a su vez proviene del aumento
\[
  \eta:\psi_*({\sh{L}'}^\bullet)
  \longrightarrow\sh{M}^{\bullet\bullet}.
\]

Para la segunda sucesión espectral, por otra parte, existen homomorfismos
de borde
\[
  {}''E_2^{i0}
  =\HH^i\!\left(
      \HH_\mathrm{I}^0
        \bigl(\Gamma(X,\sh{M}^{\bullet\bullet})\bigr)
    \right)
  \longrightarrow
  \HH^i\bigl(\Gamma(X,\sh{M}^{\bullet\bullet})\bigr),
  \tag{12.1.7.5}
  \label{0.12.1.7.5}
\]
que provienen, por \sref{0.11.3.4}, al pasar a la cohomología, del
homomorfismo de complejos
\[
  Z_\mathrm{I}^0
    \bigl(\Gamma(X,\sh{M}^{\bullet\bullet})\bigr)
  \longrightarrow\Gamma(X,\sh{M}^{\bullet\bullet}).
\]
Como $\psi_*$ es exacto por la izquierda, la sucesión
\[
  0\longrightarrow\psi_*(\sh{F}')
  \longrightarrow\psi_*({\sh{L}'}^0)
  \longrightarrow\psi_*({\sh{L}'}^1)
\]
es exacta.
Por la definición de una resolución de Cartan--Eilenberg
\sref{0.11.4.2}, podemos tomar, por tanto,
\[
  B_\mathrm{I}^0(\sh{M}^{\bullet\bullet})=0,
  \qquad
  Z_\mathrm{I}^0(\sh{M}^{\bullet\bullet})=\sh{L}^\bullet.
\]
Como el diagrama
\[
  \xymatrix@C=50pt@R=24pt{
    \sh{L}^0\ar[r]^{i^0}
      & \sh{M}^{00}\\
    \psi_*(\sh{F}')\ar[u]^\varepsilon
      \ar[r]_{\varepsilon'}\ar[ur]^{\varepsilon''}
      & \psi_*({\sh{L}'}^0)\ar[u]_{\eta^0}
  }
\]
es conmutativo, la inyección de complejos
$i:\sh{L}^\bullet\to\sh{M}^{\bullet\bullet}$ es compatible con los
aumentos $\varepsilon$ y $\varepsilon''$.
Existen, pues, dos homomorfismos de complejos de $\sh{L}^\bullet$
en $\sh{M}^{\bullet\bullet}$,
\[
  \xymatrix@C=52pt@R=22pt{
    \sh{L}^\bullet\ar[r]^i\ar[dr]_{v^\flat}
      & \sh{M}^{\bullet\bullet}\\
    & \psi_*({\sh{L}'}^\bullet)\ar[u]_\eta
  }
\]
\oldpage[0\textsubscript{III}]{57}
compatibles con los aumentos $\varepsilon$ y
$\varepsilon''$.
Como $\sh{L}^\bullet$ es una resolución inyectiva y
$\sh{M}^{\bullet\bullet}$ está constituido por haces inyectivos, se sigue
de (M, V, 1.1(a)) que estos dos homomorfismos son homotópicos.
Lo mismo vale, por tanto, para los dos homomorfismos correspondientes
\[
  \Gamma(X,\sh{L}^\bullet)
  \longrightarrow\Gamma(X,\sh{M}^{\bullet\bullet}),
\]
y, al pasar a la cohomología, dan por consiguiente el mismo
homomorfismo.
En otras palabras, hemos demostrado que el homomorfismo de borde
\sref{0.12.1.7.5}, escrito
\[
  \HH^p(X,\psi_*(\sh{F}'))
  \longrightarrow\HH^p\bigl(\Gamma(X,\sh{M}^{\bullet\bullet})\bigr),
\]
es el compuesto de \sref{0.12.1.7.1} y
\sref{0.12.1.7.3}, estando este último escrito
\[
  \HH^p(X',\sh{F}')
  \longrightarrow\HH^p\bigl(\Gamma(X,\sh{M}^{\bullet\bullet})\bigr)
\]
y habiéndose visto que es un isomorfismo.
Esto demuestra la aserción.
\end{env}

\subsection{Imágenes directas superiores}
\label{subsection:0.12.2}

\begin{env}[12.2.1]
\label{0.12.2.1}
Sean $(X,\sh{O}_X)$ y $(Y,\sh{O}_Y)$ dos espacios anillados, y sea
$f=(\psi,\theta)$ un morfismo de $X$ en $Y$.  Define el funtor imagen
directa
\[
  f_*:\cat{C}(X)\longrightarrow\cat{C}(Y),
\]
que es además idéntico a la restricción a $\cat{C}(X)$ del
funtor $\psi_*$ definido sobre la categoría de haces de grupos abelianos
en $X$.
Este último funtor es aditivo y exacto por la izquierda; y como todo haz de
grupos abelianos sobre $X$ es isomorfo a un subhaz de un haz inyectivo
de grupos abelianos, se definen los funtores derivados por la derecha
\[
  \sh{F}\longmapsto\RR^p\psi_*(\sh{F})
\]
de $\psi_*$.  Los $\RR^p\psi_*(\sh{F})$ son haces de grupos abelianos
sobre $Y$, y los $\RR^p\psi_*$ forman un funtor cohomológico universal
(T, 2.3).

Además, $\RR^p\psi_*(\sh{F})$ es el haz asociado al
prehaz
\[
  V\longmapsto\HH^p(f^{-1}(V),\sh{F})
\]
(T, 3.7.2).
Supongamos ahora que $\sh{F}$ es un $\sh{O}_X$-módulo.
Entonces $\HH^p(f^{-1}(V),\sh{F})$ está dotado naturalmente de la estructura
de $\Gamma(f^{-1}(V),\sh{O}_X)$-módulo, y por tanto de la de
$\Gamma(V,\psi_*(\sh{O}_X))$-módulo; el homomorfismo
\[
  \theta:\sh{O}_Y\longrightarrow\psi_*(\sh{O}_X)
\]
lo dota entonces de la estructura de
$\Gamma(V,\sh{O}_Y)$-módulo.
Para las estructuras así definidas, la restricción de un subconjunto abierto
$V$ a un subconjunto abierto $V'\subset V$ define un dihomomorfismo.  Esto
define, por tanto, sobre todo $\RR^p\psi_*(\sh{F})$ la estructura de un
$\sh{O}_Y$-módulo.  Denotamos este $\sh{O}_Y$-módulo por
$\RR^p f_*(\sh{F})$; así, $\RR^p f_*$ queda definido como un funtor
aditivo de $\cat{C}(X)$ en $\cat{C}(Y)$.

Además, los $\RR^p f_*$ forman un $\delta$-funtor.  En efecto, si
\[
  0\longrightarrow\sh{F}'\longrightarrow\sh{F}
  \longrightarrow\sh{F}''\longrightarrow0
\]
es una sucesión exacta de $\sh{O}_X$-módulos, la descripción precedente
de los $\RR^p\psi_*$ y de la estructura de $\sh{O}_Y$-módulo sobre
$\RR^p\psi_*(\sh{F})$ muestra inmediatamente que el homomorfismo de conexión
\[
  \partial:\RR^p\psi_*(\sh{F}'')
  \longrightarrow\RR^{p+1}\psi_*(\sh{F}')
\]
es un homomorfismo de $\sh{O}_Y$-módulos.
Finalmente, los $\RR^p f_*$ se identifican con los funtores derivados por la
derecha de $f_*$.  Todo $\sh{O}_X$-módulo admite una resolución inyectiva por
$\sh{O}_X$-módulos, y tal resolución está constituida por haces flácidos
de grupos abelianos (G, II, 7.1); puede, por tanto, utilizarse para calcular
los $\RR^p\psi_*(\sh{F})$, pues $\RR^n\psi_*(\sh{G})=0$ para $n\geq1$
y todo haz flácido $\sh{G}$ (T, 2.4.1, Observación~3, y el
 corolario de la Proposición~\EGAUnavailableHyperref{0.3.3.2}{3.3.2}).
Se sigue que los $\RR^p f_*$ forman un funtor cohomológico universal
de $\cat{C}(X)$ en $\cat{C}(Y)$ (T, 2.3).
\end{env}

\begin{env}[12.2.2]
\label{0.12.2.2}
Sean $\sh{F}$ y $\sh{G}$ dos $\sh{O}_X$-módulos.
Con la notación de \sref{0.12.2.1}, para todo subconjunto abierto $V$ de
$Y$ existe el homomorfismo de producto cup \sref{0.12.1.2.1}
\[
  \begin{aligned}
  &\HH^p(f^{-1}(V),\sh{F})
    \otimes_{\Gamma(f^{-1}(V),\sh{O}_X)}
    \HH^q(f^{-1}(V),\sh{G})\\
  &\hspace{55mm}\longrightarrow
    \HH^{p+q}\bigl(f^{-1}(V),
      \sh{F}\otimes_{\sh{O}_X}\sh{G}\bigr).
  \end{aligned}
\]
\oldpage[0\textsubscript{III}]{58}
Se sigue inmediatamente de la definición del producto cup
(G, II, 6.6) que estos homomorfismos conmutan con la restricción de
$V$ a un subespacio abierto $V'$ de $V$.
Por otra parte, $\theta$ da un homomorfismo de anillos
\[
  \Gamma(V,\sh{O}_Y)
  \longrightarrow\Gamma(V,\psi_*(\sh{O}_X))
  =\Gamma(f^{-1}(V),\sh{O}_X),
\]
y, por tanto, un homomorfismo canónico de productos tensoriales
\[
  \begin{aligned}
  &\HH^p(f^{-1}(V),\sh{F})
    \otimes_{\Gamma(V,\sh{O}_Y)}
    \HH^q(f^{-1}(V),\sh{G})\\
  &\quad\longrightarrow
    \HH^p(f^{-1}(V),\sh{F})
    \otimes_{\Gamma(f^{-1}(V),\sh{O}_X)}
    \HH^q(f^{-1}(V),\sh{G}),
  \end{aligned}
\]
que es asimismo compatible con la restricción de $V$ a $V'$.
Por composición se obtiene así un homomorfismo de
$\Gamma(V,\sh{O}_Y)$-módulos.  Para los haces asociados a los
prehaces en cuestión, define un homomorfismo canónico,
funtorial en $\sh{F}$ y $\sh{G}$,
\[
  \RR^p f_*(\sh{F})\otimes_{\sh{O}_Y}\RR^q f_*(\sh{G})
  \longrightarrow
  \RR^{p+q}f_*
    \bigl(\sh{F}\otimes_{\sh{O}_X}\sh{G}\bigr).
  \tag{12.2.2.1}
  \label{0.12.2.2.1}
\]
Para $p=q=0$, este homomorfismo se reduce a
\sref[0\textsubscript{I}]{0.4.2.2.1}.
\end{env}

\begin{proposition}[12.2.3]
\label{0.12.2.3}
Para todo $\sh{O}_X$-módulo $\sh{F}$ y todo $\sh{O}_Y$-módulo
$\sh{L}$ localmente libre de rango finito, existen isomorfismos canónicos
funtoriales
\[
  \RR^p f_*(\sh{F})\otimes_{\sh{O}_Y}\sh{L}
  \simto
  \RR^p f_*
    \bigl(\sh{F}\otimes_{\sh{O}_X}f^*(\sh{L})\bigr).
  \tag{12.2.3.1}
  \label{0.12.2.3.1}
\]
\end{proposition}

\begin{proof}
El homomorfismo \sref{0.12.2.3.1} se obtiene componiendo el
siguiente caso particular de \sref{0.12.2.2.1},
\[
  \RR^p f_*(\sh{F})
    \otimes_{\sh{O}_Y}f_*\bigl(f^*(\sh{L})\bigr)
  \longrightarrow
  \RR^p f_*
    \bigl(\sh{F}\otimes_{\sh{O}_X}f^*(\sh{L})\bigr),
  \tag{12.2.3.2}
  \label{0.12.2.3.2}
\]
con el homomorfismo del primer miembro de
\sref{0.12.2.3.1} al de \sref{0.12.2.3.2} inducido por el
homomorfismo canónico \sref[0\textsubscript{I}]{0.4.4.3.2}.
Para comprobar que \sref{0.12.2.3.1} es un isomorfismo cuando $\sh{L}$ es
localmente libre, puede reducirse inmediatamente al caso
$\sh{L}=\sh{O}_Y$, pues la cuestión es local sobre $Y$ y los
funtores en cuestión son aditivos en $\sh{L}$.
En ese caso, en vista de la definición de \sref{0.12.2.2.1}, la
proposición se reduce a comprobar que el homomorfismo correspondiente
de prehaces es biyectivo; esto se sigue inmediatamente de
$f^*(\sh{O}_Y)=\sh{O}_X$.
\end{proof}

\begin{env}[12.2.4]
\label{0.12.2.4}
Sea $(Z,\sh{O}_Z)$ un tercer espacio anillado, y sea
$g:Y\to Z$ un morfismo de espacios anillados.
Sabemos (G, II, 7.1 y 3.1.1) que, para todo
$\sh{O}_X$-módulo inyectivo $\sh{G}$, $f_*(\sh{G})$ es un haz flácido de
grupos abelianos.  Por consiguiente, por \sref{0.12.2.1},
\[
  \RR^p g_*\bigl(f_*(\sh{G})\bigr)=0
  \qquad(p>0).
\]
Se sigue (T, 2.4.1) que la sucesión espectral de Leray para funtores
compuestos se aplica al compuesto $g_*\circ f_*$.  Existe una
sucesión espectral birregular cuyo término límite es el funtor
$\RR^\bullet h_*$, donde $h=g\circ f$, y cuyo término $E_2$ es
\[
  E_2^{pq}=\RR^p g_*\bigl(\RR^q f_*(\sh{F})\bigr).
  \tag{12.2.4.1}
  \label{0.12.2.4.1}
\]
\end{env}

\begin{env}[12.2.5]
\label{0.12.2.5}
Bajo las condiciones de \sref{0.12.2.4}, definiremos directamente
homomorfismos canónicos de $\sh{O}_Z$-módulos
\[
  \RR^n g_*\bigl(f_*(\sh{F})\bigr)
  \longrightarrow\RR^n h_*(\sh{F})
  \tag{12.2.5.1}
  \label{0.12.2.5.1}
\]
y
\[
  \RR^n h_*(\sh{F})
  \longrightarrow g_*\bigl(\RR^n f_*(\sh{F})\bigr),
  \tag{12.2.5.2}
  \label{0.12.2.5.2}
\]
\oldpage[0\textsubscript{III}]{59}
que pueden identificarse con los ``homomorfismos de borde'' de la sucesión
espectral de Leray (cf.\ \sref{0.12.1.7}).
Basta trabajar con los prehaces a los que están asociados los haces de
imágenes directas superiores \sref{0.12.2.1}.
Para ello, sea $W$ un subconjunto abierto cualquiera de $Z$, y consideremos
su imagen inversa $g^{-1}(W)$ en $Y$.  Existe un
dihomomorfismo canónico
\[
  \begin{aligned}
  \HH^n\bigl(g^{-1}(W),f_*(\sh{F})\bigr)
  &\longrightarrow
  \HH^n\bigl(f^{-1}(g^{-1}(W)),
    f^*(f_*(\sh{F}))\bigr),
  \end{aligned}
  \tag{12.2.5.3}
  \label{0.12.2.5.3}
\]
siendo los anillos correspondientes
$\Gamma(g^{-1}(W),\sh{O}_Y)$ y
$\Gamma(h^{-1}(W),\sh{O}_X)$.
Por otra parte, el homomorfismo canónico
\sref[0\textsubscript{I}]{0.4.4.3.3} da, por funtorialidad,
homomorfismos canónicos
\[
  \HH^n\bigl(h^{-1}(W),f^*(f_*(\sh{F}))\bigr)
  \longrightarrow
  \HH^n\bigl(h^{-1}(W),\sh{F}\bigr),
  \tag{12.2.5.4}
  \label{0.12.2.5.4}
\]
que son homomorfismos de
$\Gamma(h^{-1}(W),\sh{O}_X)$-módulos.
Teniendo en cuenta el homomorfismo de anillos
\[
  \Gamma(W,\sh{O}_Z)
  \longrightarrow\Gamma(g^{-1}(W),\sh{O}_Y),
\]
vemos que, componiendo \sref{0.12.2.5.3} con
\sref{0.12.2.5.4}, se obtiene un homomorfismo de prehaces y, por tanto, el
homomorfismo de haces \sref{0.12.2.5.1}.

La definición de \sref{0.12.2.5.2} es aún más sencilla.
Por definición, $\RR^n h_*(\sh{F})$ está asociado al prehaz
\[
  W\longmapsto
  \HH^n\bigl(f^{-1}(g^{-1}(W)),\sh{F}\bigr),
\]
mientras que $\RR^n f_*(\sh{F})$ está asociado al prehaz
\[
  V\longmapsto\HH^n(f^{-1}(V),\sh{F}).
\]
Existe, por tanto, un homomorfismo canónico
\[
  \HH^n\bigl(f^{-1}(g^{-1}(W)),\sh{F}\bigr)
  \longrightarrow
  \Gamma\bigl(g^{-1}(W),\RR^n f_*(\sh{F})\bigr),
\]
y es inmediato que estos homomorfismos definen un homomorfismo de
prehaces, que a su vez define \sref{0.12.2.5.2}.
\end{env}

\begin{env}[12.2.6]
\label{0.12.2.6}
Bajo las hipótesis de \sref{0.12.2.4}, sean $\sh{F}$, $\sh{G}$
y $\sh{K}$ tres $\sh{O}_X$-módulos, y sea
\[
  u:\sh{F}\otimes_{\sh{O}_X}\sh{G}\longrightarrow\sh{K}
\]
un $\sh{O}_X$-homomorfismo.
Entonces los diagramas siguientes son conmutativos:
\[
  \xymatrix@C=10pt@R=24pt{
    \RR^p g_*(f_*(\sh{F}))
      \otimes_{\sh{O}_Z}\RR^q g_*(f_*(\sh{G}))
      \ar[r]\ar[d]
      & \RR^{p+q}g_*(f_*(\sh{K}))\ar[d]\\
    \RR^p h_*(\sh{F})
      \otimes_{\sh{O}_Z}\RR^q h_*(\sh{G})
      \ar[r]
      & \RR^{p+q}h_*(\sh{K}),
  }
  \tag{12.2.6.1}
  \label{0.12.2.6.1}
\]
y
\[
  \xymatrix@C=10pt@R=24pt{
    \RR^p h_*(\sh{F})
      \otimes_{\sh{O}_Z}\RR^q h_*(\sh{G})
      \ar[r]\ar[d]
      & \RR^{p+q}h_*(\sh{K})\ar[d]\\
    g_*(\RR^p f_*(\sh{F}))
      \otimes_{\sh{O}_Z}g_*(\RR^q f_*(\sh{G}))
      \ar[r]
      & g_*(\RR^{p+q}f_*(\sh{K})).
  }
  \tag{12.2.6.2}
  \label{0.12.2.6.2}
\]
\oldpage[0\textsubscript{III}]{60}
Aquí las flechas horizontales provienen de \sref{0.12.2.2.1} (la última
combinada con \sref[0\textsubscript{I}]{0.4.2.2.1}), y las flechas
verticales de \sref{0.12.2.5.1} y \sref{0.12.2.5.2}, respectivamente.

En efecto, basta verificar la aserción para los homomorfismos de prehaces
correspondientes.  Volviendo a las definiciones de
\sref{0.12.2.2} y \sref{0.12.2.5}, se reduce inmediatamente, para
\sref{0.12.2.6.1}, a los diagramas conmutativos
\sref{0.12.1.5.2}; la verificación es aún más sencilla para
\sref{0.12.2.6.2}.
\end{env}

\subsection{Suplementos sobre los funtores Ext de haces}
\label{subsection:0.12.3}

\begin{env}[12.3.1]
\label{0.12.3.1}
Sea $(X,\sh{O}_X)$ un espacio anillado.
No recordaremos la definición ni las propiedades principales de los
bifuntores
\[
  \Ext^p_{\sh{O}_X}(X;\sh{F},\sh{G}),
  \qquad
  \shExt^p_{\sh{O}_X}(\sh{F},\sh{G}),
\]
el primero con valores en la categoría de
$\Gamma(X,\sh{O}_X)$-módulos y el segundo en la categoría de
$\sh{O}_X$-módulos, ni la sucesión espectral birregular
$E(\sh{F},\sh{G})$ que los relaciona (T, 4.2 y G, II, 7.3).
\end{env}

\begin{env}[12.3.2]
\label{0.12.3.2}
Al igual que en (M, XIV, 1), se define la noción de una
\emph{extensión} de un $\sh{O}_X$-módulo $\sh{F}$ por un
$\sh{O}_X$-módulo $\sh{G}$ y la ley de composición sobre las clases de
extensiones equivalentes: los argumentos utilizados para los módulos se
trasladan manifiestamente a una categoría abeliana arbitraria.
La segunda demostración de (M, XIV, 1.1), que solo utiliza la existencia de
inmersiones en objetos inyectivos, sigue siendo, por tanto, válida en la
categoría de $\sh{O}_X$-módulos.  Muestra que
\[
  \Ext^1_{\sh{O}_X}(X;\sh{F},\sh{G})
\]
se identifica canónicamente con el grupo abeliano de clases de
extensiones de $\sh{F}$ por $\sh{G}$.
\end{env}

\begin{proposition}[12.3.3]
\label{0.12.3.3}
Sea $(X,\sh{O}_X)$ un espacio anillado cuyo haz de anillos
$\sh{O}_X$ es coherente.
Entonces, para todo par de $\sh{O}_X$-módulos coherentes $\sh{F}$ y
$\sh{G}$ y todo $p\geq0$,
\[
  \shExt^p_{\sh{O}_X}(\sh{F},\sh{G})
\]
es un $\sh{O}_X$-módulo coherente.
\end{proposition}

\begin{proof}
Los $\shExt^p_{\sh{O}_X}(\sh{F},\sh{G})$ forman un funtor cohomológico
contravariante en $\sh{F}$.
Como $\sh{F}$ es coherente, para todo $p$ y todo punto $x\in X$
existen un entorno abierto $U$ de $x$ en $X$ y una sucesión exacta
de $(\sh{O}_X|U)$-módulos
\[
  0\longrightarrow\sh{R}\longrightarrow\sh{L}_{p-1}
  \longrightarrow\cdots\longrightarrow\sh{L}_0
  \longrightarrow\sh{F}|U\longrightarrow0,
\]
donde cada $\sh{L}_i$ ($0\leq i\leq p-1$) es isomorfo a
$\sh{O}_X^{n_i}|U$, y $\sh{R}$ es coherente.
Esto se sigue por inducción sobre $p$ de
\sref[0\textsubscript{I}]{0.5.3.2} y
\sref[0\textsubscript{I}]{0.5.3.4}, pues $\sh{O}_X$ es coherente.

Además, para $p\geq1$,
\[
  \shExt^p_{\sh{O}_X|U}(\sh{L}|U,\sh{G}|U)=0
\]
para todo $\sh{O}_X$-módulo $\sh{L}$ tal que $\sh{L}|U$ sea
isomorfo a $\sh{O}_X^n|U$ (T, 4.2.3).
El argumento de (M, V, 7.2) se aplica, por tanto, al funtor cohomológico
contravariante
\[
  \sh{F}\longmapsto
  \shHom_{\sh{O}_X|U}(\sh{F}|U,\sh{G}|U)
\]
y da una sucesión exacta
\[
  \begin{aligned}
  \shHom_{\sh{O}_X|U}(\sh{L}_{p-1},\sh{G}|U)
  &\longrightarrow
  \shHom_{\sh{O}_X|U}(\sh{R},\sh{G}|U)\\
  &\longrightarrow
  \shExt^p_{\sh{O}_X|U}(\sh{F}|U,\sh{G}|U)
  \longrightarrow0.
  \end{aligned}
\]
Los dos primeros términos son $(\sh{O}_X|U)$-módulos coherentes
\sref[0\textsubscript{I}]{0.5.3.5}; por tanto, también lo es el tercero
\sref[0\textsubscript{I}]{0.5.3.4}.
\end{proof}

\oldpage[0\textsubscript{III}]{61}
\begin{proposition}[12.3.4]
\label{0.12.3.4}
Sea $f:X\to Y$ un morfismo plano de espacios anillados, y sean
$\sh{F}$ y $\sh{G}$ dos $\sh{O}_Y$-módulos.
\begin{enumerate}
  \item[\rm(i)] Existe un homomorfismo de bifuntores cohomológicos
  \[
    f^*\bigl(\shExt^\bullet_{\sh{O}_Y}(\sh{F},\sh{G})\bigr)
    \longrightarrow
    \shExt^\bullet_{\sh{O}_X}
      \bigl(f^*(\sh{F}),f^*(\sh{G})\bigr),
    \tag{12.3.4.1}
    \label{0.12.3.4.1}
  \]
  que en grado $0$ se reduce al homomorfismo canónico
  \sref[0\textsubscript{I}]{0.4.4.6}.

  \item[\rm(ii)] Existe un morfismo canónico de sucesiones espectrales
  \[
    E(\sh{F},\sh{G})
    \longrightarrow E\bigl(f^*(\sh{F}),f^*(\sh{G})\bigr),
    \tag{12.3.4.2}
    \label{0.12.3.4.2}
  \]
  cuya aplicación sobre los términos $E_2$ viene dada por los homomorfismos
  \[
    \begin{aligned}
    \HH^p\bigl(Y,\shExt^q_{\sh{O}_Y}(\sh{F},\sh{G})\bigr)
    \longrightarrow
    \HH^p\bigl(X,\shExt^q_{\sh{O}_X}
      (f^*(\sh{F}),f^*(\sh{G}))\bigr),
    \end{aligned}
    \tag{12.3.4.3}
    \label{0.12.3.4.3}
  \]
  deducidos de \sref{0.12.3.4.1} y \sref{0.12.1.3.1}.
\end{enumerate}
\end{proposition}

\begin{proof}
\begin{enumerate}
  \item[\rm(i)]
  Como $f^*$ es exacto sobre la categoría de $\sh{O}_Y$-módulos
  \sref[0\textsubscript{I}]{0.6.7.2}, los funtores
  \[
    \sh{G}\longmapsto
      f^*\bigl(\shHom_{\sh{O}_Y}(\sh{F},\sh{G})\bigr),
    \qquad
    \sh{G}\longmapsto
      \shHom_{\sh{O}_X}\bigl(f^*(\sh{F}),f^*(\sh{G})\bigr)
  \]
  son exactos por la izquierda.  El homomorfismo canónico
  \sref[0\textsubscript{I}]{0.4.4.6} induce, por tanto, canónicamente un
  homomorfismo entre sus funtores derivados.

  Para calcular estos funtores derivados, tomemos una resolución inyectiva
  $\sh{L}^\bullet=(\sh{L}^i)$ de $\sh{G}$.  Se tienen entonces
  homomorfismos
  \[
    \begin{aligned}
    \sh{H}^p\!\left(
      f^*\bigl(\shHom_{\sh{O}_Y}
        (\sh{F},\sh{L}^\bullet)\bigr)\right)
    \longrightarrow
    \sh{H}^p\!\left(
      \shHom_{\sh{O}_X}
        (f^*(\sh{F}),f^*(\sh{L}^\bullet))\right)
    \end{aligned}
  \]
  entre los haces de cohomología de complejos.
  Por la exactitud de $f^*$, el primer miembro es
  \[
    f^*\bigl(\shExt^p_{\sh{O}_Y}(\sh{F},\sh{G})\bigr).
  \]
  La exactitud de $f^*$ implica también que
  $f^*(\sh{L}^\bullet)$ es una resolución de $f^*(\sh{G})$.
  Si ${\sh{L}'}^\bullet=({\sh{L}'}^i)$ es una resolución inyectiva
  de $f^*(\sh{G})$ en la categoría de $\sh{O}_X$-módulos, existe un
  homomorfismo de complejos
  \[
    f^*(\sh{L}^\bullet)\longrightarrow{\sh{L}'}^\bullet,
  \]
  determinado salvo homotopía, y por tanto un homomorfismo bien definido en
  cohomología.  Su compuesto con el homomorfismo precedente es
  \sref{0.12.3.4.1}.

  \item[\rm(ii)]
  Con la notación precedente, existe un homomorfismo de complejos de
  $\sh{O}_X$-módulos
  \[
    f^*\bigl(\shHom_{\sh{O}_Y}(\sh{F},\sh{L}^\bullet)\bigr)
    \longrightarrow
    \shHom_{\sh{O}_X}
      \bigl(f^*(\sh{F}),{\sh{L}'}^\bullet\bigr).
  \]
  Sea $\sh{M}^{\bullet\bullet}$ una resolución inyectiva de
  Cartan--Eilenberg de
  $\shHom_{\sh{O}_Y}(\sh{F},\sh{L}^\bullet)$ en la categoría de
  $\sh{O}_Y$-módulos.  Como $f^*$ es exacto,
  $f^*(\sh{M}^{\bullet\bullet})$ es una resolución de Cartan--Eilenberg de
  $f^*(\shHom_{\sh{O}_Y}(\sh{F},\sh{L}^\bullet))$.
  Si ${\sh{M}'}^{\bullet\bullet}$ es una resolución inyectiva de
  Cartan--Eilenberg de
  $\shHom_{\sh{O}_X}(f^*(\sh{F}),{\sh{L}'}^\bullet)$, existe
  \sref{0.11.4.2} un homomorfismo, determinado salvo homotopía,
  \[
    f^*(\sh{M}^{\bullet\bullet})
    \longrightarrow{\sh{M}'}^{\bullet\bullet},
  \]
  compatible con el homomorfismo precedente; en otras palabras, existe
  un $f$-morfismo
  $\sh{M}^{\bullet\bullet}\to{\sh{M}'}^{\bullet\bullet}$ de
  de bicomplejos de haces.
  Induce un dihomomorfismo
  \[
    \Gamma(Y,\sh{M}^{\bullet\bullet})
    \longrightarrow\Gamma(X,{\sh{M}'}^{\bullet\bullet})
  \]
  de bicomplejos de módulos, determinado salvo homotopía, y por tanto
  un morfismo bien definido de sucesiones espectrales
  \sref{0.11.3.2}.  Este es precisamente el morfismo deseado
  \sref{0.12.3.4.2}, y la descripción
  \sref{0.12.3.4.3} se sigue inmediatamente de las definiciones.
\end{enumerate}
\end{proof}

\begin{proposition}[12.3.5]
\label{0.12.3.5}
Bajo las hipótesis de \sref{0.12.3.4}, supongamos además que el
haz de anillos $\sh{O}_Y$ es coherente.
Entonces, para todo $\sh{O}_Y$-módulo coherente $\sh{F}$, los
homomorfismos canónicos \sref{0.12.3.4.1} son biyectivos.
\end{proposition}

\begin{proof}
\oldpage[0\textsubscript{III}]{62}
Como la cuestión es local sobre $Y$, podemos suponer que existe una
sucesión exacta
\[
  0\longrightarrow\sh{R}\longrightarrow\sh{O}_Y^n
  \longrightarrow\sh{F}\longrightarrow0.
\]
Entonces $\sh{R}$ es también un $\sh{O}_Y$-módulo coherente
\sref[0\textsubscript{I}]{0.5.3.4}.
Para demostrar que los homomorfismos
\[
  f^*\bigl(\shExt^p_{\sh{O}_Y}(\sh{F},\sh{G})\bigr)
  \longrightarrow
  \shExt^p_{\sh{O}_X}
    \bigl(f^*(\sh{F}),f^*(\sh{G})\bigr)
\]
son biyectivos, razonamos por inducción sobre $p$.
Para $p=0$, la aserción es
\sref[0\textsubscript{I}]{0.6.7.6.1}.
Existe un diagrama conmutativo con filas exactas
\[
{\tiny
  \xymatrix@C=2pt@R=24pt{
  f^*\!\left(\shExt^{p-1}_{\sh{O}_Y}(\sh{O}_Y^n,\sh{G})\right)
    \ar[r]\ar[d]
  & f^*\!\left(\shExt^{p-1}_{\sh{O}_Y}(\sh{R},\sh{G})\right)
    \ar[r]^\partial\ar[d]
  & f^*\!\left(\shExt^p_{\sh{O}_Y}(\sh{F},\sh{G})\right)
    \ar[r]\ar[d]
  & f^*\!\left(\shExt^p_{\sh{O}_Y}(\sh{O}_Y^n,\sh{G})\right)
    \ar[d]\\
  \shExt^{p-1}_{\sh{O}_X}(\sh{O}_X^n,f^*(\sh{G}))
    \ar[r]
  & \shExt^{p-1}_{\sh{O}_X}(f^*(\sh{R}),f^*(\sh{G}))
    \ar[r]^\partial
  & \shExt^p_{\sh{O}_X}(f^*(\sh{F}),f^*(\sh{G}))
    \ar[r]
  & \shExt^p_{\sh{O}_X}(\sh{O}_X^n,f^*(\sh{G})).
  }}
\]
Aquí hemos utilizado $f^*(\sh{O}_Y)=\sh{O}_X$; como $f^*$ es exacto,
ambas filas son exactas.
Además,
\[
  \shExt^p_{\sh{O}_Y}(\sh{O}_Y^n,\sh{G})=0,
  \qquad
  \shExt^p_{\sh{O}_X}(\sh{O}_X^n,f^*(\sh{G}))=0
  \qquad(p>0)
\]
(T, 4.2.3).
Por la hipótesis de inducción, las dos primeras flechas verticales del
diagrama precedente son isomorfismos, mientras que los términos de la
derecha son nulos.  Se sigue que
\[
  f^*\bigl(\shExt^p_{\sh{O}_Y}(\sh{F},\sh{G})\bigr)
  \longrightarrow
  \shExt^p_{\sh{O}_X}
    \bigl(f^*(\sh{F}),f^*(\sh{G})\bigr)
\]
es un isomorfismo.
\end{proof}

\subsection{Hipercohomología del funtor imagen directa}
\label{subsection:0.12.4}

\begin{env}[12.4.1]
\label{0.12.4.1}
Sean $(X,\sh{O}_X)$ y $(Y,\sh{O}_Y)$ dos espacios anillados, y sea
$f:X\longrightarrow Y$ un morfismo de espacios anillados.
Se puede formar la hipercohomología de $f_*$ respecto de un complejo
arbitrario de $\sh{O}_X$-módulos
$\sh{K}^\bullet=(\sh{K}^i)_{i\in\bb{Z}}$
\sref{0.11.4.4}, pues los límites inductivos filtrantes existen y son exactos
en la categoría abeliana de $\sh{O}_Y$-módulos (T, 3.1.1).
Los $\sh{O}_Y$-módulos de hipercohomología
$\RR^p f_*(\sh{K}^\bullet)$ se denotarán también
$\sh{H}^p(f,\sh{K}^\bullet)$, o simplemente
$\sh{H}^p_f(\sh{K}^\bullet)$.
Recordemos que $\sh{H}^\bullet(f,\sh{K}^\bullet)$ es la cohomología del
bicomplejo de $\sh{O}_Y$-módulos $f_*(\sh{L}^{\bullet\bullet})$,
donde $\sh{L}^{\bullet\bullet}$ es una resolución inyectiva de
Cartan--Eilenberg de $\sh{K}^\bullet$ en la categoría de
$\sh{O}_X$-módulos.
Es el término límite de dos sucesiones espectrales cuyos términos $E_2$
vienen dados por
\begin{equation}
\label{0.12.4.1.1}
\tag{12.4.1.1}
  {}'E_2^{pq}
  =\sh{H}^p\bigl(\sh{H}^q(f,\sh{K}^\bullet)\bigr)
\end{equation}
y
\begin{equation}
\label{0.12.4.1.2}
\tag{12.4.1.2}
  {}''E_2^{pq}
  =\sh{H}^p\bigl(f,\sh{H}^q(\sh{K}^\bullet)\bigr)
  \quad
  \bigl(=\RR^p f_*(\sh{H}^q(\sh{K}^\bullet))\bigr).
\end{equation}
En estas fórmulas hemos adoptado la notación general $T(A^\bullet)$
para el transformado de un complejo por un funtor \sref{0.11.2.1}, y, para
un $\sh{O}_X$-módulo $\sh{F}$, escribimos
$\sh{H}^p(f,\sh{F})$ en lugar de $\RR^p f_*(\sh{F})$.
Recordemos también que la primera sucesión espectral es siempre regular.
Ambas sucesiones espectrales son birregulares cuando $\sh{K}^\bullet$ está
acotado inferiormente,
\oldpage[0\textsubscript{III}]{63}
o cuando existe un entero $m$ tal que todo $\sh{O}_X$-módulo
admite una resolución flácida de longitud $\leq m$ \sref{0.11.4.4}.
\end{env}

\begin{env}[12.4.2]
\label{0.12.4.2}
Denotaremos análogamente por $\bb{H}^\bullet(X,\sh{K}^\bullet)$ la
hipercohomología del funtor $\Gamma$ respecto de un complejo
$\sh{K}^\bullet$ de $\sh{O}_X$-módulos.
Así, los $\bb{H}^p(X,\sh{K}^\bullet)$ son
$\Gamma(X,\sh{O}_X)$-módulos.
Además, $\bb{H}^\bullet(X,\sh{K}^\bullet)$ puede considerarse como un
caso particular de $\sh{H}^\bullet(f,\sh{K}^\bullet)$, donde $f$ es un
morfismo de $(X,\sh{O}_X)$ a un espacio anillado constituido por un punto
dotado del anillo $\Gamma(X,\sh{O}_X)$.

Para todo subconjunto abierto $V$ de $X$, escribiremos
$\bb{H}^\bullet(V,\sh{K}^\bullet)$ en lugar de
$\bb{H}^\bullet(V,\sh{K}^\bullet|V)$.
\end{env}

\begin{proposition}[12.4.3]
\label{0.12.4.3}
Para todo entero $p\in\bb{Z}$, el $\sh{O}_Y$-módulo
$\sh{H}^p(f,\sh{K}^\bullet)$ es canónicamente isomorfo al haz
asociado al prehaz
\[
  U\longmapsto
  \bb{H}^p\bigl(f^{-1}(U),\sh{K}^\bullet\bigr)
\]
sobre $Y$.
\end{proposition}

\begin{proof}
Con la notación de \sref{0.12.4.1}, el haz de cohomología
$\sh{H}^p(f_*(\sh{L}^{\bullet\bullet}))$ está asociado al
prehaz
\[
  U\longmapsto
  \HH^p\bigl(\Gamma(U,f_*(\sh{L}^{\bullet\bullet}))\bigr)
  =
  \HH^p\bigl(\Gamma(f^{-1}(U),\sh{L}^{\bullet\bullet})\bigr).
\]
Pero está claro que
$\sh{L}^{\bullet\bullet}|f^{-1}(U)$ es una resolución inyectiva de
Cartan--Eilenberg de $\sh{K}^\bullet|f^{-1}(U)$ (T, 3.1.3).
Por consiguiente,
\[
  \HH^p\bigl(\Gamma(f^{-1}(U),\sh{L}^{\bullet\bullet})\bigr)
  =
  \bb{H}^p\bigl(f^{-1}(U),\sh{K}^\bullet\bigr)
\]
por definición.
\end{proof}

\begin{proposition}[12.4.4]
\label{0.12.4.4}
La hipercohomología $\sh{H}^\bullet(f,\sh{K}^\bullet)$ es un
funtor cohomológico de $\sh{K}^\bullet$ en cada uno de los casos
siguientes:
\begin{enumerate}
  \item[\rm{a)}] $\sh{K}^\bullet$ varía en la categoría de complejos
    acotados inferiormente;
  \item[\rm{b)}] existe un entero $m$ tal que todo
    $\sh{O}_X$-módulo admite una resolución flácida de longitud $\leq m$;
  \item[\rm{c)}] $X$ es un espacio noetheriano.
\end{enumerate}
\end{proposition}

\begin{proof}
Los casos~a) y~b) son casos particulares de \sref{0.11.5.4}.
El caso~c), por otra parte, se sigue de \sref{0.11.5.2}, pues en
este caso el funtor $f_*$ conmuta con los límites inductivos
(G, II, 3.10.1).
\end{proof}

\begin{env}[12.4.5]
\label{0.12.4.5}
Sea ahora $\mathfrak{U}=(U_\alpha)$ un recubrimiento abierto de $X$.
Para todo complejo de \emph{prehaces}
$\sh{K}^\bullet=(\sh{K}^j)$ sobre $X$, consideremos el bicomplejo
$\check{\mathrm{C}}^\bullet(\mathfrak{U},\sh{K}^\bullet)$ cuya
componente de bigrado $(i,j)$ es
$\check{\mathrm{C}}^i(\mathfrak{U},\sh{K}^j)$, el grupo de
$i$-cocadenas alternadas del nervio de $\mathfrak{U}$ con valores
en $\sh{K}^j$ (G, II, 5.1).
Llamamos a la cohomología de este bicomplejo la
\emph{hipercohomología del recubrimiento $\mathfrak{U}$ con
coeficientes en $\sh{K}^\bullet$}, y la denotamos por
\[
  \bb{H}^\bullet(\mathfrak{U},\sh{K}^\bullet)
  =
  \HH^\bullet
  \bigl(\check{\mathrm{C}}^\bullet
    (\mathfrak{U},\sh{K}^\bullet)\bigr).
\]
La sucesión espectral de Leray de un recubrimiento
(T, 3.8.1 y G, II, 5.9.1) se generaliza a la hipercohomología como sigue.
\end{env}

\begin{proposition}[12.4.6]
\label{0.12.4.6}
Sea $\sh{K}^\bullet=(\sh{K}^i)$ un complejo de
$\sh{O}_X$-módulos.
Existe un funtor espectral regular de $\sh{K}^\bullet$ con
término límite $\bb{H}^\bullet(X,\sh{K}^\bullet)$ y término $E_2$
\begin{equation}
\label{0.12.4.6.1}
\tag{12.4.6.1}
  E_2^{pq}
  =
  \CHH^p\bigl(\mathfrak{U},h^q(\sh{K}^\bullet)\bigr),
\end{equation}
donde $h^q(\sh{K}^\bullet)$ denota el complejo de prehaces
\[
  V\longmapsto\HH^q(V,\sh{K}^\bullet)
\]
sobre $X$.
Esta sucesión espectral es birregular si $\sh{K}^\bullet$ está acotado
inferiormente.
\end{proposition}

\begin{proof}
Sea $\sh{L}^{\bullet\bullet}$ una resolución inyectiva de
Cartan--Eilenberg de $\sh{K}^\bullet$, y consideremos el tricomplejo
\[
  \check{\mathrm{C}}^\bullet
  (\mathfrak{U},\sh{L}^{\bullet\bullet})
  =
  \bigl(\check{\mathrm{C}}^i
    (\mathfrak{U},\sh{L}^{jk})\bigr).
\]
Consideremos primero este tricomplejo como un bicomplejo en los grados $i$ y
$j+k$.
Como $i$ solo toma valores no negativos, la segunda sucesión espectral
de este bicomplejo es regular \sref{0.11.3.3} y degenerada, porque
\[
  \CHH^q(\mathfrak{U},\sh{L}^{jk})=0
  \qquad(q>0),
\]
siendo los $\sh{O}_X$-módulos $\sh{L}^{jk}$ haces flácidos
\oldpage[0\textsubscript{III}]{64}
(G, II, 5.2.3).
Se sigue que existe un isomorfismo canónico
\[
  \HH^n\bigl(
    \check{\mathrm{C}}^\bullet
      (\mathfrak{U},\sh{L}^{\bullet\bullet})
  \bigr)
  \isoto
  \HH^n\bigl(\Gamma(X,\sh{L}^{\bullet\bullet})\bigr)
\]
\sref{0.11.1.6}, en virtud de (G, II, 5.2.2), y por tanto, por la
definición \sref{0.12.4.2}, un isomorfismo
\[
  \HH^n\bigl(
    \check{\mathrm{C}}^\bullet
      (\mathfrak{U},\sh{L}^{\bullet\bullet})
  \bigr)
  \isoto
  \bb{H}^n(X,\sh{K}^\bullet).
\]

Consideremos ahora
$\check{\mathrm{C}}^\bullet
(\mathfrak{U},\sh{L}^{\bullet\bullet})$
como un bicomplejo en los grados $i+j$ y $k$.
Como $k$ solo toma valores no negativos, la primera sucesión espectral
de este bicomplejo es siempre regular.
Es birregular si $\sh{L}^{jk}=0$ para $j<j_0$, es decir, cuando
$\sh{K}^\bullet$ está acotado inferiormente \sref{0.11.3.3}.
Esta es la sucesión espectral deseada.
En efecto, para todo $j$, $\sh{L}^{j,\bullet}$ es una resolución
inyectiva de $\sh{K}^j$; por consiguiente, cuando $i$ varía,
\[
  \HH^q\bigl(
    \check{\mathrm{C}}^i
      (\mathfrak{U},\sh{L}^{j,\bullet})
  \bigr)
\]
es precisamente el complejo de cocadenas
$\check{\mathrm{C}}^i(\mathfrak{U},h^q(\sh{K}^j))$.
Esto demuestra la proposición.
\end{proof}

\begin{corollary}[12.4.7]
\label{0.12.4.7}
Si, para todo símplex $\sigma$ del nervio de $\mathfrak{U}$ y
todo entero $i$, se tiene
\[
  \HH^q(U_\sigma,\sh{K}^i)=0
  \qquad(q>0),
\]
entonces existe un isomorfismo canónico
\begin{equation}
\label{0.12.4.7.1}
\tag{12.4.7.1}
  \bb{H}^\bullet(\mathfrak{U},\sh{K}^\bullet)
  \isoto
  \bb{H}^\bullet(X,\sh{K}^\bullet).
\end{equation}
\end{corollary}

\begin{proof}
La hipótesis implica que
$\check{\mathrm{C}}^\bullet
(\mathfrak{U},h^q(\sh{K}^j))=0$ para $q>0$, y por tanto que
$E_2^{pq}=0$ para $q>0$.
Como la sucesión espectral \sref{0.12.4.6.1} es degenerada y
regular, la conclusión se sigue de la definición
\sref{0.12.4.5} de
$\bb{H}^\bullet(\mathfrak{U},\sh{K}^\bullet)$ y de
\sref{0.11.1.6}.
\end{proof}

\begin{env}[12.4.8]
\label{0.12.4.8}
Sea $(X',\sh{O}_{X'})$ un segundo espacio anillado, y sea
$f=(\psi,\theta)$ un morfismo de $X'$ en $X$.
Por el mismo método que en \sref{0.12.1.3} y \sref{0.12.1.4}, se
define un dihomomorfismo para la hipercohomología de un complejo
$\sh{K}^\bullet$ de $\sh{O}_X$-módulos:
\begin{equation}
\label{0.12.4.8.1}
\tag{12.4.8.1}
  \bb{H}^p(X,\sh{K}^\bullet)
  \longrightarrow
  \bb{H}^p(X',f^*(\sh{K}^\bullet)).
\end{equation}
Partiendo de una resolución inyectiva de Cartan--Eilenberg
$\sh{L}^{\bullet\bullet}$ de $\sh{K}^\bullet$, observemos que, como
$\psi^*$ es exacto, $\psi^*(\sh{L}^{\bullet\bullet})$ es una
resolución de Cartan--Eilenberg de $\psi^*(\sh{K}^\bullet)$ en la
categoría de $\psi^*(\sh{O}_X)$-módulos.
Existe, por tanto, un morfismo
\[
  \psi^*(\sh{L}^{\bullet\bullet})
  \longrightarrow
  {\sh{L}'}^{\bullet\bullet},
\]
donde ${\sh{L}'}^{\bullet\bullet}$ es una resolución inyectiva de
Cartan--Eilenberg de $\psi^*(\sh{K}^\bullet)$.
Induce un morfismo en cohomología
\[
  \bb{H}^\bullet(X,\sh{K}^\bullet)
  \longrightarrow
  \bb{H}^\bullet
  (X',\psi^*(\sh{K}^\bullet)).
\]
Componiéndolo con el morfismo inducido por funtorialidad a partir de
$\psi^*(\sh{K}^\bullet)\longrightarrow f^*(\sh{K}^\bullet)$ se obtiene
el morfismo deseado \sref{0.12.4.8.1}.

Partiendo de \sref{0.12.4.8.1} y \sref{0.12.4.3}, y razonando como
en \sref{0.12.2.5}, se pueden definir entonces, para dos morfismos de
espacios anillados
\[
  f:X\longrightarrow Y,
  \qquad
  g:Y\longrightarrow Z,
\]
con $h=g\circ f$, homomorfismos para la hipercohomología de un complejo
$\sh{K}^\bullet$ de $\sh{O}_X$-módulos:
\begin{equation}
\label{0.12.4.8.2}
\tag{12.4.8.2}
  \sh{H}^n(g,f_*(\sh{K}^\bullet))
  \longrightarrow
  \sh{H}^n(h,\sh{K}^\bullet),
\end{equation}
y
\begin{equation}
\label{0.12.4.8.3}
\tag{12.4.8.3}
  \sh{H}^n(h,\sh{K}^\bullet)
  \longrightarrow
  g_*\bigl(\sh{H}^n(f,\sh{K}^\bullet)\bigr).
\end{equation}
Dejamos al lector los detalles de las definiciones.
\end{env}
